\section{Reboques}

\subsection{Reboques civis}

As matrículas dos reboques civis seguem sequências próprias: de \plate{X-0} a
\plate{X-000000}, ou de \plate{XX-0} a \plate{XX-00000}, em que \plate{X} ou
\plate{XX} é o código do serviço distrital ou regional onde o reboque foi
registado (ver o anexo~\ref{sec:distritos}).

\platepair{trailer-av}{Reboque de Aveiro: \plate{AV-45152}.}%
          {trailer-l-174771}{Reboque de Lisboa: \plate{L-174771}.}

\subsection{Reboques civis de Lisboa}

Em Lisboa, o número do reboque permite estimar a data do registo:

\gen{trailer-lisbon}

\subsection{Reboques militares}

\begin{itemize}
  \item Os reboques do Exército têm matrículas com as sequências \plate{0-M},
        \plate{00-M}, \plate{000-M} e \plate{0000-M}.
  \item Os reboques da Força Aérea têm a sequência \plate{E-00-00}.
  \item Os reboques da Marinha têm a mesma sequência dos seus automóveis:
        \plate{AP-00-00}.
\end{itemize}

O número de matrícula é colocado numa chapa de formato e cores iguais às dos
automóveis, mas sem a barra amarela com a data de matrícula. Por desconhecimento
da norma, é comum verem-se chapas de reboques com a barra amarela e a data de
registo.
